\section{Discussion}
\label{sec:discussion}

\paragraph{Where the gain comes from}
In every environment the search structure accounts for most of the advantage of the learning-guided method, and the structure is not the same twice: destruction--reconstruction for identical machines (the IG beats the ILS by
$0.3$--$0.7\%$ with learning removed from both sides), its combination with a descent for unrelated machines ($4.2\%$; $2.1\%$ at equal CPU on the standard benchmark) and the local search for the flowshop. The
learned policy adds $0.01$--$0.05$ points for identical machines, a difference that no class resolves but that a pooled test confirms on held-out and fresh published instances, and about $0.1\%$ for unrelated machines;
in the flowshop it subtracts. A budget confound sits on top, and it is partly an implementation effect: $28$--$30\%$ of the advantage of a local search for identical machines is budget when the descent is vectorized or compiled ($93\%$ with a first, loop-based version), the descent is neutral for the recipe at equal CPU (it lost by $0.034\%$ with the loop version), and the agent
itself costs $1.4$--$1.7\times$ the CPU time of random selection for unrelated machines. Studies that compare at an equal number of iterations and without an offline-tuned competitor would have credited these gains to learning; Table~\ref{tab:fair_vs_naive} lists six places where the protocol changes the reading.

\begin{table}[t]
\centering
\caption{Where the budget-fair protocol changes the reading. ``Common practice'' compares at an equal number of iterations or against random selection only; ``budget-fair'' compares at equal CPU time (calibrated
sequentially) or against the configuration validated offline. $\bar d$ in \% of the makespan (negative: first method better).}
\label{tab:fair_vs_naive}
{\scriptsize\setlength{\tabcolsep}{3pt}
\begin{tabular}{@{}p{0.31\textwidth}p{0.20\textwidth}p{0.20\textwidth}p{0.17\textwidth}@{}}
\toprule
Comparison & Common practice & Budget-fair & Reading \\
\midrule
P$||C_{\max}$, IG: add a local search & $-0.150$ (equal iterations) & $-0.105$ (equal CPU, vectorized); $-0.115$ (compiled); $-0.010$ (loop) & $28$--$30\%$ budget; $93\%$ with the slow implementation \\
P$||C_{\max}$, recipe: add a local search & $-0.036$ (equal iterations) & $-0.008$ (vectorized); $-0.006$ (compiled); $+0.034$ (loop) & neutral; reverses with the slow implementation \\
R$||C_{\max}$ benchmark: IG with local search vs ILS & $-3.23$ (equal $T$) & $-2.14$ (equal CPU) & a third smaller \\
R$||C_{\max}$ $1000\times50$: Q-Learning vs random selection & $-0.35$ (equal $T$; n.s.) & $-0.19$ (equal CPU; n.s.) & nearly halved; not confirmed \\
P$||C_{\max}$, destruction size: Q-Learning vs random / vs validated & $+0.015$ (vs random) & $+0.103$ (vs validated) & online loses by seven times more \\
R$||C_{\max}$ benchmark, destruction size: Q-Learning vs random / vs validated & $+0.152$ (vs random) & $+0.360$ (vs validated) & online loses by $2.4$ times more \\
\bottomrule
\end{tabular}}
\end{table}

\paragraph{When a policy can help}
Several explanations of the small policy effect were tested. The headroom of the action space is necessary: where the per-instance oracle gains little (three destruction strategies, identical machines: $0.04\%$) the agent has
nothing to learn, and where it is larger (unrelated machines: $0.38\%$) the agent captures $30$--$38\%$ of it ($0.44$--$0.68$ of the validated gain at $1000\times50$ and on the standard benchmark), and the gain follows the headroom across cells (Spearman $0.59$ and $0.74$). It is not sufficient:
for the destruction size, with headroom that validation captures, Q-Learning is worse than random choice in all three problems, also when pre-trained or run with five times the budget, and the correlation between headroom and gain
reverses ($-0.37$). A contextual bandit, two neural Q-learners (the larger with replay and a target network, also in a joint space of $15$ actions, where the headroom is no larger: $0.11\%$) and two selectors trained offline did not change this, except that a selector trained offline on $700$ instances matches the validated configuration (Section~\ref{sec:unrelated}): learning pays when it uses data of other instances, not the first iterations of one run. The exploratory diagnostics point to the learning window (more than $90\%$ of the reward arrives in the first tenth of the run) and to the type of action, a hypothesis for pre-specified follow-up. Together, the
evidence supports a simple rule for practice: \emph{choose the configuration offline on separate instances; consider an online policy only where the actions are qualitatively different and one of them dominates}. Where qualitatively different actions are complementary and none dominates, as the four operators of the ILS, the static competitor to beat is a mixture, not the best single action, and none of three learners (tabular, contextual, a $2\times64$ DQN) beat the uniform mixture there (Section~\ref{sec:unrelated}).

\paragraph{What works}
The strongest methods of the study contain no learned component. Iterated Greedy with a validated destruction size and a start from the best constructive solution is provably optimal on $97.5\%$ of the synthetic identical-machine instances
at $64{,}000$ iterations; for unrelated machines, an IG started from the linear relaxation and polished by a restricted MILP is within $0.25$--$0.30\%$ of the best known solutions of the standard benchmark, comparable to CPLEX with two hours
($0.18$--$0.20\%$) though behind Gurobi and the exact method of the benchmark ($0.07$--$0.1\%$). The contribution is modest as a solver (LP rounding and restricted MILPs are known ideas) and is reported as such: it shows that, once the
attribution is done, the best use of the effort was in the structure and the budget, not in the agent.

\paragraph{A portable protocol}
Six elements of the protocol are inexpensive to add to other studies and changed the conclusions here: a random-selection control at matched budget separates mechanism from policy; equal CPU time, calibrated sequentially, separates
structure from budget; an offline-validated configuration is the competitor that a practitioner would use; the oracle headroom tells what any selector could gain; pooled tests and seed averaging detect small effects that per-class
single-run tests cannot; and hypotheses and decision rules with hashes, fixed before the runs, separate confirmed effects from effects generated by looking at the data (an exploratory effect shrank by more than half on held-out
instances, and a headroom law suggested by one sample was refuted prospectively). The protocol is a methodology contribution that generalises beyond scheduling to any adaptive component of a search.

\paragraph{Limits}
(i) The main agent is tabular, with one dense reward and one pre-training scheme tried, and the action spaces are small; six learners of other classes were tested (a tabular policy-gradient learner in spaces of up to $30$ actions; a linear contextual bandit and a one-hidden-layer neural Q-learner in all three problems; two instance-conditional selectors trained offline on $280$ and $700$ instances in unrelated machines only, the one with heterogeneous instances; a $2\times64$ DQN with replay and a target network, in two problems and in a joint space of $15$ actions), \todo{and deep policy-gradient learners, portfolios of hundreds of operators and tighter budgets are outside the evidence, so the results are about these learners and action spaces, not about learning in general} \todonote{to improve: test deep policy-gradient learners and portfolios of hundreds of operators}. (ii) The solvers are not state of the art: exact methods are stronger for both parallel-machine problems (about $85\%$ of families C and F of the identical-machine benchmark are solved
in about a second; $0.07$--$0.1\%$ from the best known for unrelated machines), and the heuristic is weakest on job-correlated and small instances. (iii) Most instances for unrelated machines and the flowshop are synthetic or generated; the standard benchmark was
used for one validation and the classic flowshop instances only for a replication with $40$ instances. (iv) The main comparisons use $T=1{,}000$ iterations ($T=300$ for the flowshop); the equal-CPU analyses show that conclusions about added components depend on the budget
convention and on the efficiency of the implementation (a compiled implementation of the search gives the same reading as the vectorized one; a count of evaluations agrees in sign but depends on its cost model; the tabular, bandit and policy-gradient agents for identical machines were compiled and cost $1.0$--$1.3$ times random choice per unit of work, but the neural agents and the unrelated-machines pipeline remain in Python, where the agent costs $1.4$--$1.7$ times random choice) \todonote{to improve: compile the unrelated-machines pipeline and the neural agents and check their cost ratio}, and the policy was not studied at budgets much beyond $T=5{,}000$. (v) \todo{The pre-specification is a dated, hash-stamped diary, not a public registry} \todonote{to improve: Zenodo deposit with a DOI}, and the number of analyses is large; all families and verdicts, with every non-confirmation,
are in Supplementary Table~\ref{S-tab:s_prereg_all}. (vi) The synthesis across action types was noticed after the results and is a hypothesis.
